\documentclass[12pt]{article}
\usepackage[utf8]{inputenc}
\usepackage[T1]{fontenc}
\usepackage[margin=1in]{geometry}
\usepackage{setspace}
\usepackage{xcolor}
\usepackage[colorlinks=true, urlcolor=blue!60!black, linkcolor=blue!60!black]{hyperref}
\setlength{\parindent}{2em}
\setlength{\parskip}{1.25em}
\renewcommand{\baselinestretch}{1.0}
\begin{document}
\begin{center}
{\LARGE \textit{United States v. Nixon}, 418 U.S. 683 (1974)\par}
\vspace{0.8em}
{\large Case Brief --- Four-Page Edition\par}
\end{center}
\vspace{1em}

\section{Memory Jogger}

In a unanimous 8--0 decision arising from Watergate, the Court held that while presidential communications enjoy a constitutionally rooted presumptive privilege, a President's generalized interest in confidentiality must yield to a demonstrated, specific need for relevant and admissible evidence in a pending federal criminal trial (\href{https://supreme.justia.com/cases/federal/us/418/683/}{418 U.S. 683}).

\section{Detailed Case Facts}

After the June 1972 Watergate break-in, Alexander Butterfield revealed in July 1973 that Nixon had secretly taped Oval Office conversations. Following the ``Saturday Night Massacre'' firing of Special Prosecutor Archibald Cox, Acting Attorney General Robert Bork appointed Leon Jaworski under a Justice Department regulation giving the Special Prosecutor power to contest executive privilege and protecting him from removal except for ``extraordinary improprieties.''

In March 1974 a grand jury indicted seven former aides and officials, including John Mitchell, H. R. Haldeman, and John Ehrlichman, and named Nixon an unindicted coconspirator. On the Special Prosecutor's motion under Rule 17(c), the district court subpoenaed tapes and papers of 64 identified conversations. Nixon released edited transcripts of some conversations and moved to quash, claiming absolute privilege and nonjusticiability.

\section{Procedural History}

Chief Judge John Sirica denied the motion and ordered the originals produced for \textit{in camera} inspection. Nixon appealed to the D.C. Circuit; the Special Prosecutor sought certiorari before judgment, which the Court granted, along with Nixon's cross-petition on the coconspirator naming. After argument on July 8, 1974, the Court ruled on July 24. Following release of the ``smoking gun'' tape, Nixon resigned effective August 9, 1974.

\section{Judicial Votes}

Unanimous, \textbf{8--0}.
\begin{itemize}
\item \textbf{Majority:} Chief Justice Burger, joined by Justices Douglas, Brennan, Stewart, White, Marshall, Blackmun, and Powell; three Nixon appointees joined the judgment against him.
\item \textbf{Concurrences and dissents:} None.
\item \textbf{Not participating:} Justice Rehnquist, who had served in Nixon's Justice Department.
\end{itemize}

\section{Holding}

The order compelling \textit{in camera} production was affirmed; Nixon's cross-petition was dismissed as improvidently granted. (1) An order denying a President's motion to quash is immediately appealable. (2) The dispute is justiciable because the Executive is bound by its own regulation under \textit{Accardi}. (3) Rule 17(c) requires relevancy, admissibility, and specificity. (4) Presidential communications are presumptively privileged, but absent military, diplomatic, or national security secrets a generalized claim of confidentiality yields to the specific need for evidence in a criminal trial. Footnote 19 reserved civil discovery, congressional demands, and state secrets.

\section{Analysis of Opinions}

\textbf{Jurisdiction.} Under the \textit{Perlman} exception, forcing a President into contempt to obtain review would be ``unseemly.''

\textbf{Justiciability.} Rejecting Nixon's ``intra-branch'' political question argument, Burger looked past a ``surface inquiry'': the case was a federal prosecution, and the unrevoked regulation supplied the ``concrete adverseness'' Article III requires.

\textbf{Rule 17(c).} Under \textit{Bowman Dairy} and \textit{Iozia}, the subpoena met each requirement even under the ``particularly meticulous'' review owed a President under \textit{United States v. Burr}.

\textbf{Executive privilege.} Reaffirming \textit{Marbury}, Burger held that the judiciary interprets the scope of privilege. He recognized for the first time a constitutionally based ``presumptive privilege'' implied from Article II, but, quoting Justice Jackson in \textit{Youngstown}, rejected an absolute privilege because the Constitution mandates ``separateness but interdependence, autonomy but reciprocity.'' Weighing confidentiality against the principle that ``the public has a right to every man's evidence'' and the Fifth and Sixth Amendments, he found occasional disclosure would not chill advisers, while absolute privilege would ``gravely impair'' the courts. The district court was to review with scrupulous care and return excised material under seal.

\section{Examples: Future Applications}

\begin{enumerate}
\item \textbf{Same-Side: Bribery Prosecution.} A special counsel subpoenas recordings of three identified Oval Office meetings about a bribe, and the President asserts only generalized confidentiality. The subpoena is enforceable for \textit{in camera} review.
\item \textbf{Same-Side: State Grand Jury.} A state prosecutor subpoenas a sitting President's pre-presidency financial records from his accountants. There is no absolute immunity from criminal process (as in \textit{Trump v. Vance}, 2020).
\item \textbf{Opposite-Side: Military War Plans.} A defendant subpoenas National Security Council briefings on deployments and intelligence sources. \textit{Nixon} promised ``utmost deference'' to military and diplomatic secrets, so the privilege would likely prevail.
\item \textbf{Opposite-Side: Civil Fishing Expedition.} A civil plaintiff seeks two years of White House emails. The request lacks specificity, and \textit{Nixon} did not extend its balance to civil litigation.
\item \textbf{Fence-Sitter: Congressional Subpoena.} A House committee seeks memoranda between the President and advisers. \textit{Nixon} left this unresolved: Congress lacks the Sixth Amendment imperative but has implied investigative power, so courts favor accommodation.
\end{enumerate}

\section{Critique}

\textbf{Originalist and textualist views.} Raoul Berger and Saikrishna Prakash question a constitutional privilege the text never mentions, contrasting Article I's Speech or Debate Clause. Unitary-executive theorists argue that a subordinate cannot sue the President and that a revocable regulation could not create an Article III controversy.

\textbf{Progressive and doctrinal views.} William Van Alstyne observed that in seeking unanimity the Court defeated Nixon while handing future presidents a constitutional privilege later used against congressional oversight. Critics fault the vague balancing test, and commentators note tension between \textit{Nixon} and \textit{Trump v. United States} (2024).

\section{Key Quotations}

\begin{quote}
``A President and those who assist him must be free to explore alternatives in the process of shaping policies and making decisions, and to do so in a way many would be unwilling to express except privately. These are the considerations justifying a presumptive privilege for Presidential communications. The privilege is fundamental to the operation of Government, and inextricably rooted in the separation of powers under the Constitution.''
\end{quote}
\noindent --- Chief Justice Burger, Opinion of the Court, 418 U.S.\ at 708.

\begin{quote}
``We conclude that, when the ground for asserting privilege as to subpoenaed materials sought for use in a criminal trial is based only on the generalized interest in confidentiality, it cannot prevail over the fundamental demands of due process of law in the fair administration of criminal justice. The generalized assertion of privilege must yield to the demonstrated, specific need for evidence in a pending criminal trial.''
\end{quote}
\noindent --- Chief Justice Burger, Opinion of the Court, 418 U.S.\ at 713.

\begin{quote}
``However, neither the doctrine of separation of powers nor the need for confidentiality of high-level communications, without more, can sustain an absolute, unqualified Presidential privilege of immunity from judicial process under all circumstances.''
\end{quote}
\noindent --- Chief Justice Burger, Opinion of the Court, 418 U.S.\ at 706.

\end{document}
